\section{Method}

\subsection{Worker, Quant-H0, and Evolvable Harness}

Let $M$ be a frozen base model and $H_i$ the harness at candidate version $i$.
The resulting Worker Agent is
\begin{equation}
    A_i = \operatorname{Agent}(M,H_i), \qquad M_i=M_0.
\end{equation}
The harness may contain prompts, executable tools, memory, middleware,
validators, routing, and research workflow. One execution of task $q$ produces
a trajectory $\tau_{i,q}$ and artifact $a_{i,q}$.

The main method starts from Quant-H0-S6-Core, a newly frozen baseline containing the
Quant Research Worker identity, one shell tool, and six concise observable
stage checkpoints. The system prompt makes the generalization object explicit:
the Worker must complete the current public task using reusable research
capability rather than a hidden-evaluator or task-answer patch. Exact
task-specific rules remain valid when the current public instruction, supplied
public data, or a predeclared public reference states them.

The same minimum provenance rule is injected by the trusted Worker runtime, so
an evolved candidate cannot remove it by editing its own prompt. This
instruction-level boundary cannot replace Candidate Information-Set Review:
the prospective main controller must determine whether the exact cumulative
effective candidate may reach a Worker.

The thin state interface is a registered skill loaded before substantive shell
work. It fixes state names, marker grammar, public-grounded inapplicability, and
evaluation-to-earlier-stage revisits, but does not prescribe a data-audit
checklist, quantitative method, perturbation, validator, or artifact algorithm.
Each stage emits only a short public summary---not hidden chain of
thought---for mandate, evidence, representation, operation, evaluation, and
artifact completion. Evaluation may explicitly revisit an earlier stage, and
an inapplicable stage must be marked with a public reason rather than silently
skipped. A trusted post-run parser stores the marker ledger outside the task
submission and does not equate marker presence with stage correctness.
This is a human-authored observation scaffold shared by the baseline and QRS;
the claimed QRS contribution begins only when state-conditioned evidence,
relation selection, component routing, or intervention changes the harness.
The more detailed S6-Full workflow is retained as a disclosed manual-process
baseline, not silently folded into the QRS starting state. Quant-H0-S6-Core
contains no evolved component library, candidate history, failure classifier,
or component-selection policy.

\subsection{Quant Research State Representation}

We use the general state set
\begin{equation}
    \mathcal{S}=\{S_{\mathrm{mandate}},S_{\mathrm{evidence}},
    S_{\mathrm{representation}},S_{\mathrm{operation}},
    S_{\mathrm{evaluation}},S_{\mathrm{artifact}}\}.
\end{equation}
The states denote Research Mandate \& Contract, Research Evidence \& Data,
Quantitative Representation, Research Operation, Evaluation \& Reconciliation,
and Research Artifact \& Completion. They are an open diagnostic ontology, not
a claim that research is intrinsically linear. Quant-H0-S6-Core nevertheless orders
their observable checkpoints to make histories comparable and diagnosable;
evaluation can record a revisit to evidence, representation, or operation.
Research Operation may instantiate as estimation, pricing, calibration, risk
analysis, backtesting, execution, or another operation.

For the relevant part of a failed trajectory, the Evolver reconstructs a
task-conditioned expected state $s^*_q$ and observed state
$\tilde{s}_{i,q}$, then predicts a target state $\hat{s}_{i+1,q}$. It compares
at least two plausible explanations before selecting the earliest consequential
mismatch supported by the available trajectory and artifact evidence.

\subsection{Task-Conditioned Transition Observations}

From the public task contract, the system may select a small set
$\mathcal{I}(q)$ of applicable executable relations. Each invariant returns
\begin{equation}
    v_k(q,\tau,a) \in
    \{\Pass,\Fail,\NA,\Unknown\}.
\end{equation}
Candidate families include temporal prefix consistency, estimator fit scope,
quantity relations, portfolio and execution identities, cost monotonicity,
universe perturbations, and fresh artifact replay. Applicability is
task-conditioned. An invariant is one possible observation of a Research State
transition; it is not the complete state representation or a replacement for
the official benchmark verifier.

\subsection{Research-State-Guided Component Search}

An evolution proposal records
\begin{equation}
    h_i = (s^*_q,\tilde{s}_{i,q},c_i,\hat{s}_{i+1,q},o_i,
    \widehat{\Delta r}_i),
\end{equation}
where $c_i$ is a selected or synthesized harness component, $o_i$ is a concrete
trajectory, component-use, invariant, artifact, or official observation of the
predicted transition, and $\widehat{\Delta r}_i$ is the predicted official
outcome or protected behavior. The mutation surface remains open; a Research
State mismatch narrows the search question without forcing a fixed component
mapping.

Component experience retains positive, negative, inactive, and unstable
episodes. Retrieval is conditioned on the Research State, attempted component,
observed transition, and official outcome rather than only task identity or a
broad language failure class.

When runtime evidence does not locate a state mismatch, the Evolver may return
calibrated insufficiency using the same State Card and authorized history
interface. A proposal that returns \texttt{ACT} must also enumerate every
decision-changing claim that would reach the Worker and identify the prompt,
tool, descriptor, middleware, memory, or configuration surface that carries it.

\subsection{Candidate Information-Set Review}

The prospective QRS main Evolver receives only public and answer-free evidence.
Historical answer-rich development is retained as a threat model rather than
as main-method permission: retrospective audit showed that a hidden evaluator
predicate can be semantically projected into a Worker-visible prompt or tool
even when the original diagnostic file is absent. Raw-file isolation alone is
therefore insufficient.

Under the prospective main contract, every admitted changed candidate enters a mandatory candidate-focused
Candidate Information-Set Review before fresh Worker execution. A trusted
controller constructs the review package from the complete cumulative
Quant-H0-S6-Core-to-effective-candidate Worker-visible material, the candidate's claim inventory,
public task contracts, and any independent references frozen before proposal.
Optimize-only material may be supplied only to trace a forbidden origin; it
cannot positively support a claim. The Reviewer checks both support and
coverage. Each decision-changing claim must be entailed by a public contract or
a predeclared independent reference, and every changed Worker-visible exposure
must be declared. Overall \texttt{PASS} and complete coverage permit Worker
evaluation. \texttt{REJECT} or \texttt{INCONCLUSIVE} returns
\texttt{HOLD\_FOR\_REFINE} or stops the lineage without a Worker call.

The Candidate Information-Set Reviewer is method-agnostic and has no
authority to edit, evaluate, retain, or promote a harness. Its \texttt{PASS}
means only that the candidate is eligible for blind evaluation. A separate
callable Quant Research Reviewer may remain an optional search aid that returns
competing explanations or a discriminating observation; it is not the
Candidate Information-Set Reviewer and is not required by the main treatment.

\subsection{Intervention Verdicts}

We keep six observations separate:

\begin{enumerate}
    \item \textbf{Information-set eligibility}: the complete candidate receives
    Candidate Information-Set Review \texttt{PASS}.
    \item \textbf{Component activation or prompt exposure}: the fresh Worker
    reaches a callable component, or receives a prompt treatment by
    construction; exposure alone does not establish behavioral uptake.
    \item \textbf{Research-State correction}: the predeclared transition
    observation changes as predicted.
    \item \textbf{Official target improvement}: the unchanged benchmark verifier
    reports a property or task gain.
    \item \textbf{Stable promotion}: an independent repeat reproduces the
    resolved-property footprint and a matched answer-free protection task is
    aggregate- and property-set safe.
    \item \textbf{Sealed performance}: a frozen final incumbent is evaluated
    once on tasks that never entered search or selection.
\end{enumerate}

State correction can provide useful harness evidence before a complete solve,
but it is not reported as benchmark improvement. An official gain without the
predicted transition remains performance evidence with unresolved mechanism
attribution. We maintain separate official and search parents: binary-reward
Pareto improvement updates the official incumbent, while a non-regressing
candidate with observed component activation and a supported predeclared state
transition may advance only the search parent. \Measured{R4 exercised the
mechanics of dual-parent refinement, independent repeat, protection, and
controller promotion.} \Boundary{The later provenance audit found that R4's
Worker-visible SVI rule projected a diagnostic-only predicate. Its official
scores and controller transitions remain measured, but R4 is contaminated
development evidence and not a valid scientific promotion or a clean example
of the proposed information boundary.}

\Proposed{The main unit is an independent QRS cumulative campaign rather than a
set of disconnected task-specific winners. Each campaign starts from a fresh
Quant-H0-S6-Core, maintains one official incumbent, and traverses a frozen family
order under fixed evidence, mutation, and evaluation budgets. A separately
executed QRS-no-State ablation uses the same lifecycle without State-Card
construction or state-conditioned retrieval/routing; it is not a runtime mode
in the released method. AHE is reproduced in its own pinned implementation and
compared externally on common frozen sealed outcomes and resource use. Each
final incumbent is frozen before a no-feedback evaluation whose outcomes never
return to proposal, retrieval, task selection, or promotion. QuantCodeEval
breadth is a useful later cross-benchmark test but does not block the first
QFBench QRS campaign.}
